\hypertarget{appendix-aa405b7c3903}{}
\section*{Borrow contracts and pool epochs}
\section{Borrow contracts and pool epochs}\label{borrow-contracts-and-pool-epochs}

\textbf{Curriculum depth:} Introduced · \textbf{First seen:} C008 · \textbf{Latest development:} C008

\subsection{Definition and mechanism}\label{definition-and-mechanism}

An epoch versions a borrow-validity interval. A checked borrow compares its saved epoch to the live pool counter before using the cached element. The pool counter must outlive the borrow; counters alone do not supply ownership, synchronization, race freedom or wrap immunity.

\subsection{Prerequisites}\label{prerequisites}

\begin{itemize}
\tightlist
\item
  \href{}{Borrow invalidation and logical identity (05-Engine-I, appendix-f9b8b9188c7f)} --- Invalidating mutations define the deadlines a saved epoch or deferred borrow must respect.
\end{itemize}

\subsection{Used by}\label{used-by}

Later sources will deepen this same node; no empty dependent nodes are created.

\subsection{Recall and next depth}\label{recall-and-next-depth}

List every mutation that advances the epoch and identify check-then-use races. Pending: production wrappers, wrap policy and concurrent lifetime pinning. Written coverage does not establish learner mastery.

\subsection{Source and navigation}\label{source-and-navigation}

\href{Continuous\%20Notes/Computing\%20Foundations\%20-\%20From\%20Bits\%20to\%20C++\%20and\%20Systems\%20APIs\#When\%20a\%20valid\%20borrow\%20stops\%20naming\%20the\%20intended\%20thing}{Teaching explanation and examples} · \hyperref[appendix-fd090d982a15]{Exact coverage} · \href{}{Dependency Map (02-Foundations-II, appendix-4ad47bdd269a)}
